\chapter{Arquitectura Metodologica del Sistema}

\section{3.1 Vision general del pipeline experimental}
El sistema se implementa como pipeline secuencial de adquisicion, transformacion y validacion de evidencia. Cada etapa publica artefactos que alimentan la etapa siguiente.

\section{3.2 Estructura modular del software}
Se explicita el reparto funcional entre codigo fuente, herramientas de validacion y scripts de ejecucion de campanas.

\section{3.3 Contrato de artefactos y trazabilidad}
Se define correspondencia entre archivos de entrada, salida y metadatos de ejecucion para asegurar reproducibilidad.

\section{3.4 Definicion de plan y su representacion}
Se describe la representacion del plan de asignacion como estructura consumible por simulador y runtime.

\section{3.5 Criterios de validez y riesgos metodologicos iniciales}
Se documentan condiciones de calidad estadistica, consistencia topologica y amenazas de validez para etapas tempranas.
